% ============================================================================
% deepdive_part23.tex -- Parts II and III of the RISP mathematical
% deep dive. This is a FRAGMENT: it assumes the master preamble defines
% theorem/proposition/lemma/corollary/definition/example/remark (numbered by
% section), the tcolorbox environments intuition/keypoint/warning, and the
% macros \R \E \Var \KL \TV \nm \argmax \argmin.
% All labels are prefixed sec:dd2- / eq:dd2- / tab:dd2-.
% ============================================================================

% ============================================================================
% PART II: THE ENVIRONMENT
% ============================================================================
\part{The Environment: Regimes, Episodes, and Dormancy}
\label{sec:dd2-part-env}

The mechanism documented in Part~III is an answer to a question posed
entirely by the environment: regimes that arrive in persistent blocks
(Section~\ref{sec:dd2-regimes}), episodes as the true statistical unit
(Section~\ref{sec:dd2-episodes}), and the combinatorics of dormancy, during
which a capacity-limited learner is quietly destroyed
(Section~\ref{sec:dd2-dormancy}). Every design choice in \nm{} is a response
to a quantitative fact established in this part.

\section{Regime-Switching Markets}
\label{sec:dd2-regimes}

\subsection{The Regime Process}

\begin{definition}[Block-structured regime process]
\label{def:dd2-regime-process}
Fix a finite regime set $\mathcal{R} = \{1, \ldots, R\}$. A
\textbf{block-structured (semi-Markov) regime process} is a sequence of pairs
\[
(\rho_1, \ell_1),\ (\rho_2, \ell_2),\ (\rho_3, \ell_3),\ \ldots
\qquad \rho_b \in \mathcal{R},\ \ell_b \in \{1, 2, \ldots\},
\]
where $\rho_b$ is the regime of the $b$-th \emph{block} and $\ell_b$ its
duration in days, with $\rho_{b+1} \neq \rho_b$ (blocks are maximal). The
day-level regime label $r_t$ is the piecewise-constant expansion: writing
$T_b = \sum_{b' < b} \ell_{b'}$ for the start of block $b$,
\begin{equation}
\label{eq:dd2-semimarkov}
r_t = \rho_b \qquad \text{for } T_b < t \le T_b + \ell_b .
\end{equation}
The process is semi-Markov when $\rho_{b+1}$ depends on the past only through
$\rho_b$ (a transition kernel on $\mathcal{R}$) while $\ell_b$ is drawn from a
regime-specific duration distribution $G_{\rho_b}$ on $\{1,2,\ldots\}$.
\end{definition}

Hamilton's Markov-switching model \citep{hamilton1989new}---the canonical
formalization of regimes in econometrics---is the special case in which $r_t$
is a first-order Markov chain: block durations are then forced to be
geometric, $\Pr(\ell = k) = P_{rr}^{\,k-1}(1 - P_{rr})$ with mean
$1/(1 - P_{rr})$. The semi-Markov form frees the duration law $G_r$, which
matters because the experimental stream uses bounded, non-geometric blocks
(25--60 days common, 15--30 days crisis) and because real regime durations
are empirically far from geometric \citep{ang2012regime, guidolin2007asset}.
Nothing in Part~III depends on the duration law beyond persistence itself.

The concrete instantiation used throughout: $R = 4$ regimes
$\mathcal{R} = \{\textsf{trend}, \textsf{range}, \textsf{highvol},
\textsf{crisis}\}$; common regimes arrive in blocks of 25--60 days;
\textsf{crisis} arrives in short blocks of 15--30 days and only on every
\emph{fourth} cycle through the common regimes, so consecutive crisis
episodes are separated by roughly 250--450 days of dormancy; horizon
$T \approx 2370$ days ($\approx 9.5$ trading years); each day carries
features $X_t \in \R^{n \times d}$ ($n = 20$ assets, $d = 20$ features per
asset) and realized returns $y_t \in \R^{n}$.

\subsection{Why Block Persistence (Not I.I.D.\ Regime Draws) Is the Whole Game}

It is tempting to model regime-switching as $r_t \stackrel{iid}{\sim} \pi$
for some distribution $\pi$ on $\mathcal{R}$. That model destroys both of the
phenomena this project is about.

\begin{proposition}[Persistence creates local stationarity and dormancy]
\label{prop:dd2-persistence}
Let $r_t$ be a regime process over $\mathcal{R}$.
\begin{enumerate}
  \item \textbf{(I.i.d.\ draws have no usable present.)} If
        $r_t \stackrel{iid}{\sim} \pi$, the conditional law of $r_{t+1}$
        given the entire past equals $\pi$: knowing today's regime says
        nothing about tomorrow's, every day is a fresh mixture draw, and the
        optimal single-day strategy is mixture-optimal, not any
        regime-specialist.
  \item \textbf{(I.i.d.\ draws have no dormancy.)} Under i.i.d.\ draws,
        $\Pr(\text{regime } r \text{ dormant} > D \text{ days})
        = (1 - \pi(r))^{D}$, exponentially small in $D$. Under the block
        process, dormancy is the sum of intervening block durations and is
        typically hundreds of days when $r$ sits on a long cycle---as
        \textsf{crisis} does here ($250$--$450$ days by construction).
\end{enumerate}
\end{proposition}

\begin{proof}
(1) is immediate from independence. (2): with
$\pi(\textsf{crisis}) \approx 1/13$ (one short crisis block per $\approx 12$
common blocks), an i.i.d.\ stream would produce a 250-day crisis dormancy
with probability $(1 - 1/13)^{250} \approx e^{-19.2} \approx 5 \times
10^{-9}$---essentially never---whereas the block schedule produces one every
cycle. The block construction makes long dormancy the rule rather than a
large-deviation event.
\end{proof}

\begin{intuition}[title={Persistence is what makes specialization purchasable}]
Persistence does two jobs at once. \emph{Within} a block, the world is locally
stationary for 25--60 days, long enough for a specialist to be trained,
evaluated, and trusted---this is what makes a regime a \emph{niche} worth
owning. \emph{Between} blocks of the same regime, the world is silent about
that regime for weeks to years---this is what makes memory a problem. Under
i.i.d.\ draws neither holds: there is nothing to specialize in (job one fails)
and nothing is ever dormant (job two never arises). Every result in this
document lives in the gap between those two timescales.
\end{intuition}

\subsection{The Causal Labelers}
\label{sec:dd2-labelers}

All eleven arms of the experiment receive the regime label $\hat r_t$ from a
shared \textbf{causal labeler}: a deterministic function of information
available \emph{strictly before} the decision at day $t$ is made. Two labelers
are used, deliberately spanning a simple/structured spectrum, so that no
conclusion depends on one labeling artifact.

\paragraph{Labeler L1: rolling-percentile volatility band $\times$ trend sign.}
At day $t$, using returns $y_1, \ldots, y_{t-1}$ only:
\begin{enumerate}
  \item \textbf{Trailing volatility.} Realized volatility
        $v_t = \operatorname{sd}\bigl(\bar y_{t-w_v}, \ldots, \bar y_{t-1}\bigr)$
        over a trailing window of $w_v$ days, $\bar y_s$ the cross-sectional
        mean return on day $s$.
  \item \textbf{Rolling percentile thresholds.} The empirical
        $q^{\text{high}}$- and $q^{\text{crisis}}$-quantiles
        $\tau^{\text{high}}_t, \tau^{\text{crisis}}_t$ of the trailing
        history $\{v_s : s \le t-1\}$---quantiles of \emph{past}
        volatilities only, re-estimated daily from the filtration at $t-1$,
        never from the full sample.
  \item \textbf{Trend statistic.} A trailing trend signal $m_t$ (sign and
        magnitude of the mean of $\bar y_{t-w_m}, \ldots, \bar y_{t-1}$, or
        a fast-minus-slow moving-average gap), again from days $\le t-1$.
  \item \textbf{Label.}
        \[
        \hat r_t =
        \begin{cases}
          \textsf{crisis} & v_t \ge \tau^{\text{crisis}}_t, \\
          \textsf{highvol} & \tau^{\text{high}}_t \le v_t < \tau^{\text{crisis}}_t, \\
          \textsf{trend} & v_t < \tau^{\text{high}}_t \text{ and } |m_t| \ge \tau^{m}, \\
          \textsf{range} & v_t < \tau^{\text{high}}_t \text{ and } |m_t| < \tau^{m}.
        \end{cases}
        \]
\end{enumerate}

\paragraph{Labeler L2: forward-only filtered two-state volatility model
$\times$ trend sign.}
L2 replaces the percentile band with a probabilistic state model:
\begin{enumerate}
  \item Maintain a two-state Markov-switching volatility model (low/high-vol
        states with state-dependent return variance), parameters estimated
        on an expanding window through day $t-1$ and refit periodically from
        past data only.
  \item Run the \textbf{Hamilton filter} \citep{hamilton1989new}
        \emph{forward only} for the one-step-ahead predictive probability
        $p_t = \Pr(s_t = \text{high-vol} \mid \bar y_1, \ldots, \bar y_{t-1})$.
        The Kim \emph{smoothed} probability, which conditions on the future,
        is never computed.
  \item Label \textsf{crisis} when $p_t$ exceeds a high threshold and
        trailing volatility is extreme, \textsf{highvol} when $p_t$ exceeds
        a moderate threshold, otherwise \textsf{trend}/\textsf{range} by the
        same trend statistic $m_t$ as in L1.
\end{enumerate}

\paragraph{Leakage analysis.}
Write $\mathcal{F}_{t-1} = \sigma(X_1, y_1, \ldots, X_{t-1}, y_{t-1}, X_t)$
for the information legitimately available when the day-$t$ decision is
placed (today's features are observable; today's returns are not). Both
labelers satisfy the \textbf{causality contract}
\begin{equation}
\label{eq:dd2-causal}
\hat r_t \in \mathcal{F}_{t-1}
\qquad \text{(in fact } \hat r_t \in \sigma(y_1, \ldots, y_{t-1})\text{)},
\end{equation}
because every ingredient---trailing volatility, rolling quantiles, trend
statistics, filter recursions, parameter refits---is a function of returns
through $t-1$. Three subtler channels are also closed:
\emph{(i) contemporaneous leak:} the label for day $t$ never uses $y_t$
itself (a one-day embargo between the last return used and the day labeled);
\emph{(ii) refit leak:} parameters and thresholds are re-estimated only on
data through $t-1$, so no full-sample estimate back-propagates into early
labels;
\emph{(iii) smoothing leak:} L2 uses filtered, never smoothed,
probabilities---the smoother is two-sided, and a single smoothed label can
encode information from arbitrarily far in the future.

\begin{warning}[title={Label leakage: the classic regime-paper sin}]
A large fraction of the regime-switching literature evaluates strategies
against regime labels that \emph{could not have been known at the time}. The
recurring offenders:
\begin{itemize}
  \item \textbf{Smoothed state probabilities:} using Kim-smoothed
        $\Pr(s_t \mid y_{1:T})$ as ``the regime at $t$''---the smoother
        conditions on the future and labels turning points in advance.
  \item \textbf{Full-sample thresholds:} defining \textsf{highvol} as the
        top volatility quintile \emph{of the whole sample}---the 1990
        boundary then depends on 2008.
  \item \textbf{Peak/trough dating:} NBER-style bull/bear datings require
        future data to \emph{confirm} a peak; using the dated label at the
        peak itself is hindsight.
  \item \textbf{Full-sample parameters with filtered probabilities:} the
        leak hides in $\hat\theta$, not in the recursion.
\end{itemize}
Each inflates the apparent value of regime-conditioning for the same reason
published return predictors decay out of sample \citep{mclean2016does}: the
evaluation quietly conditions on information the strategy never had.
Everything downstream---specialist training, competition windows,
pinning---consumes only the causal label
$\hat r_t \in \mathcal{F}_{t-1}$ of \eqref{eq:dd2-causal}; ``reactivation''
means the \emph{causal labeler} says so, with whatever lag and
misclassification that entails.
\end{warning}

\begin{remark}[The label is an input, not a contribution]
\nm{} takes $\hat r_t$ as given and identical across all arms: the
experiment isolates \emph{what to do with} a regime signal (allocation,
training, memory), not how to detect regimes---a separate, well-studied
problem \citep{nystrup2018dynamic, nystrup2020greedy}. A shared causal label
removes regime-detection skill as a confound.
\end{remark}

\section{Episodes as Environments}
\label{sec:dd2-episodes}

\subsection{Episodes}

\begin{definition}[Episode]
\label{def:dd2-episode}
For regime $r \in \mathcal{R}$, the \textbf{$e$-th episode of $r$}, written
$(r, e)$, is the $e$-th maximal block of consecutive days on which
$\hat r_t = r$, i.e.\ the $e$-th element of the subsequence of blocks
$\{b : \rho_b = r\}$ in Definition~\ref{def:dd2-regime-process}. We write
$\mathcal{T}_{r,e} \subset \{1, \ldots, T\}$ for its day set and
$|\mathcal{T}_{r,e}| = \ell_{r,e}$ for its length.
\end{definition}

The central modeling move---inherited from invariant-prediction thinking and
from the decision-focused learning literature
\citep{elmachtoub2022smart, yan2026invpnco}---is to treat each episode as a
distinct \emph{environment}: the regime fixes what is stable, the episode
index fixes what is incidental.

\begin{definition}[Hyper-distribution assumption]
\label{def:dd2-hyper}
Each regime $r$ carries a \textbf{hyper-distribution} $Q_r$ over environment
parameters. Episode $(r,e)$ draws its parameter
$\gamma_{r,e} \sim Q_r$ independently across $e$, and the days within the
episode are then i.i.d.\ given $(r, \gamma_{r,e})$:
\[
\gamma_{r,e} \stackrel{iid}{\sim} Q_r, \qquad
(X_t, y_t) \mid (r, \gamma_{r,e}) \stackrel{iid}{\sim} P_{r, \gamma_{r,e}}
\quad \text{for } t \in \mathcal{T}_{r,e}.
\]
The \textbf{stable component} of regime $r$ is whatever is shared by
$P_{r,\gamma}$ across all $\gamma$ in the support of $Q_r$; the
\textbf{episode-specific component} is the rest.
\end{definition}

\begin{remark}[Two honest idealizations]
\label{rem:dd2-idealizations}
Definition~\ref{def:dd2-hyper} idealizes reality in two ways.
\emph{(i) Macro memory.} Real episodes of the same regime are not independent
draws: the 2008 and 2011 crises share causal machinery and participants who
remember the previous one. Positive dependence across $e$ means the effective
number of independent episodes is \emph{smaller} than the count---scarcity
(Section~\ref{sec:dd2-scarcity}) is, if anything, understated.
\emph{(ii) Secular drift.} $Q_r$ is treated as fixed over the horizon, but
markets drift on decadal scales: predictors decay after publication
\citep{mclean2016does}, and the ``stable'' component of a regime is stable
only relative to the horizon. The synthetic generator below satisfies both
idealizations exactly, so failures of an arm on synthetic data cannot be
blamed on them; transfer to long real panels inherits both caveats.
\end{remark}

\subsection{Episode Scarcity and Its Two Mitigations}
\label{sec:dd2-scarcity}

The statistical currency of Definition~\ref{def:dd2-hyper} is the number of
\emph{episodes} $E_r$, not the number of days: any quantity that varies at the
episode level (an episode-mean loss, an episode-specific coefficient) is
estimated with variance $\propto 1/E_r$, no matter how many days each episode
contains. For common regimes $E_r$ is comfortable (dozens of 25--60-day
blocks over $T \approx 2370$ days). For \textsf{crisis} it is brutal: in
roughly 35 years of equity data one finds on the order of \textbf{5--7}
crisis episodes, and fewer in shorter panels; in the synthetic stream, the
every-fourth-cycle schedule produces a comparable handful over
$T \approx 2370$ days. Two standard mitigations trade scarcity for
dependence, and it is important to state exactly what dependence each one
buys.

\begin{definition}[Sub-episode blocks]
\label{def:dd2-subepisode}
Fix a block length $B$. The \textbf{sub-episode blocks} of episode $(r,e)$
are the contiguous segments obtained by partitioning $\mathcal{T}_{r,e}$ into
pieces of length $B$ (discarding or merging a short remainder). Each block is
treated as a pseudo-episode in episode-level statistics.
\end{definition}

Sub-episode blocks multiply the apparent environment count by
$\approx \ell_{r,e}/B$, but all blocks of one episode \emph{share the same
draw} $\gamma_{r,e}$, so block mean losses are positively correlated through
$\gamma$, and cross-block variance \emph{underestimates} cross-episode
variance: the blocks inform about within-episode noise, not about $Q_r$.
They are a variance mitigation for the stable component only; they cannot
manufacture information about episode heterogeneity the panel does not
contain.

\begin{definition}[Cross-asset episodes]
\label{def:dd2-crossasset}
When the same regime episode is observed simultaneously on $A$ assets or
markets (here $n = 20$ assets share each episode), the per-asset day records
$\{(x_{t,i}, y_{t,i}) : t \in \mathcal{T}_{r,e}\}$ for $i = 1, \ldots, A$ may
be treated as $A$ parallel realizations of the episode environment.
\end{definition}

Cross-asset episodes multiply the day count per episode by $A$, but
introduce \emph{cross-sectional dependence}: assets within a day share
common shocks (market factors, and the episode draw $\gamma_{r,e}$ itself,
common across assets in the generator below), so the effective sample size
per day lies between $1$ and $A$. Again: more data about the within-episode
law, no new draws from $Q_r$.

\begin{keypoint}
Both mitigations enlarge the denominator of \emph{within-episode} estimates
and leave the number of independent draws from $Q_r$ untouched. Whatever
must be learned \emph{about episode heterogeneity}---in particular the
variance-across-episodes term the INV objective penalizes---must come from
$E_r$ genuine episodes, and for \textsf{crisis} $E_r$ is single-digit. This
one fact sizes the episode library ($\le 6$ retained episodes suffices,
because there are never many more) and explains the crisis-pinning wrinkle
of Section~\ref{sec:dd2-competition}.
\end{keypoint}

\subsection{The Synthetic Generator: Heterogeneity Made Explicit}
\label{sec:dd2-generator}

The synthetic stream instantiates Definition~\ref{def:dd2-hyper} exactly. For
day $t$ in episode $(r,e)$ and asset $i$ with feature vector
$x_{t,i} \in \R^{20}$:
\begin{equation}
\label{eq:dd2-generator}
y_{t,i} \;=\;
\underbrace{\theta_r \cdot (x_{t,i,1},\, x_{t,i,2})}_{\text{stable, regime-level}}
\;+\;
\underbrace{\gamma_{r,e}\, x_{t,i,\mathrm{sp}}}_{\text{episode-specific}}
\;+\; b_r \;+\; \varepsilon_{t,i},
\end{equation}
with the following calibration:
\begin{itemize}
  \item \textbf{Stable signal.} $\theta_r \in \R^2$, $\|\theta_r\| = 0.012$,
        directions well-separated across regimes---so a head trained on one
        regime is wrong (not merely stale) on another.
  \item \textbf{Episode-specific signal.}
        $\gamma_{r,e} \sim \mathcal{N}\bigl(0, (0.012 \cdot \mathrm{het})^2\bigr)$,
        \emph{redrawn independently each episode}, loading on a designated
        spurious feature $x_{\mathrm{sp}}$; $\mathrm{het} \ge 0$ is the
        heterogeneity knob.
  \item \textbf{Distractors.} 16 of the 20 features carry no signal in any
        regime: overfitting bait at $n$--$d$ parity ($n = d = 20$).
  \item \textbf{Noise.} $\varepsilon_{t,i} \sim \mathcal{N}(0, \sigma_r^2)$,
        $\sigma_r = 0.02 \cdot (1,\, 0.75,\, 2,\, 3)$ for
        (\textsf{trend}, \textsf{range}, \textsf{highvol}, \textsf{crisis}):
        the rare regime is also the noisiest, by a factor of three.
\end{itemize}

\begin{example}[What $\mathrm{het} = 0$ versus $\mathrm{het} = 2$ means]
\label{ex:dd2-het}
The standard deviation of the episode-specific coefficient is
$\operatorname{sd}(\gamma_{r,e}) = 0.012 \cdot \mathrm{het}$, against a stable
coefficient magnitude of $\|\theta_r\| = 0.012$. Concretely:
\begin{itemize}
  \item $\mathbf{het = 0}$: $\gamma_{r,e} \equiv 0$. Every episode of regime
        $r$ has the \emph{same} law; the spurious feature is a 17th
        distractor; $Q_r$ is a point mass; plain ERM on pooled regime data
        is right, and a variance-across-episodes penalty has nothing to
        penalize.
  \item $\mathbf{het = 1}$: $\operatorname{sd}(\gamma) = 0.012 = \|\theta_r\|$.
        The spurious coefficient is at least as large as the stable one
        whenever $|\gamma_{r,e}| \ge 0.012$, probability
        $\Pr(|Z| \ge 1) \approx 0.317$ for $Z \sim \mathcal{N}(0,1)$: in
        roughly one episode in three, the most profitable in-episode feature
        is the one that betrays you next episode.
  \item $\mathbf{het = 2}$: $\operatorname{sd}(\gamma) = 0.024 = 2\|\theta_r\|$.
        Now $|\gamma_{r,e}| \ge \|\theta_r\|$ with probability
        $\Pr(|Z| \ge 0.5) \approx 0.617$: in the \emph{majority} of episodes
        the episode-specific signal dominates inside the episode. An
        in-episode ERM fit loads heavily on $x_{\mathrm{sp}}$---correctly,
        for that episode---then carries a coefficient of typical size
        $0.024$ with a freshly randomized sign into the next episode.
\end{itemize}
The probe window (the first 15 days after a reactivation,
Definition~\ref{def:dd2-probe}) is precisely where this bites: the learner
meets a new draw of $\gamma_{r,e}$ before any in-episode data has
accumulated, so whatever $\gamma$-chasing is baked into its head is realized
as pure decision regret.
\end{example}

\begin{intuition}[title={The generator is a forgetting--overfitting double trap}]
Equation~\eqref{eq:dd2-generator} arms two independent failure modes. The
\emph{regime} structure (separated $\theta_r$, long dormancy) punishes arms
that forget; the \emph{episode} structure ($\gamma_{r,e}$ redrawn) punishes
arms that remember too literally---memorize the last crisis episode including
its $\gamma$ and you import a coin-flip coefficient into the next one. In the
eleven arms of Section~\ref{sec:dd2-arms}, the allocation axis addresses the
first trap, the objective axis (ERM vs.\ INV) the second.
\end{intuition}

\section{Dormancy Combinatorics and Memory Models}
\label{sec:dd2-dormancy}

\subsection{Dormancy and the Probe Window}

\begin{definition}[Dormancy]
\label{def:dd2-dormancy}
The \textbf{dormancy} of regime $r$ before episode $e+1$ is the gap
\[
D_r(e) \;=\; \min \mathcal{T}_{r,e+1} \;-\; \max \mathcal{T}_{r,e} \;-\; 1,
\]
the number of days between the end of episode $(r,e)$ and the start of
$(r,e+1)$ during which other regimes are active. In the synthetic stream,
common regimes have dormancies on the order of one cycle (tens of days),
while \textsf{crisis} has $D \approx 250$--$450$ days.
\end{definition}

\begin{definition}[Probe window and forgetting deficit]
\label{def:dd2-probe}
A \textbf{reactivation} of $r$ is the first day of an episode $(r, e+1)$ with
dormancy $D_r(e) \ge 90$ days. The \textbf{probe window} is the first
$15$ days of that episode. For an arm $A$, let
$\bar\rho_{\mathrm{probe}}(A)$ denote its mean per-day decision regret
(Definition~\ref{def:dd2-regret}) over probe-window days. The
\textbf{forgetting deficit} of arm $A$ at dormancy $D$ is
\begin{equation}
\label{eq:dd2-phi}
\Phi(D; A) \;=\;
\E\bigl[\bar\rho_{\mathrm{probe}}(A) \,\big|\, \text{dormancy } D\bigr]
\;-\;
\E\bigl[\bar\rho_{\mathrm{probe}}(A^{\circ})\bigr],
\end{equation}
where $A^{\circ}$ is the \emph{retention counterfactual}: the same arm
serving the probe window with the head (and episode library) it held at the
close of the previous episode of $r$, untouched by any intervening foreign
training. $\Phi \approx 0$ means dormancy cost the arm nothing it had;
$\Phi$ large means the dormancy machinery below did its damage.
\end{definition}

The probe window is 15 days: long enough to average regime noise
($\sigma_{\textsf{crisis}} = 0.06$ per day) into a stable regret estimate,
short enough that within-episode relearning cannot mask a cold start---it
measures what you walked in with. All headline numbers in this document
($0.0303$--$0.0320$) are means of $\bar\rho_{\mathrm{probe}}$ across
reactivations and 20 seeds.

\subsection{The Foreign-Update Process}

During $r$'s dormancy, the stream keeps producing days from the $W_a$
\emph{other} active regimes, and the holder of $r$'s head may be made to
\emph{train on a foreign regime}. We abstract this as the
\textbf{foreign-update process}: each dormancy day is a foreign-update day
with probability $p$, where $p$ encodes the arm: $p = 1$ for a monolith that
trains daily on whatever is active; $p \approx$ (win rate) for an unpinned
\nm{} specialist; $p = 0$ for a pinned owner whose regime is dormant. The
damage per foreign update is what the two memory models specify.

\begin{definition}[Hard and soft memory models]
\label{def:dd2-memory-models}
Each specialist has capacity $K < R$: at most $K$ regime heads (linear SPO+
heads, Section~\ref{sec:dd2-specialist}), each with an episode buffer of up to
$6$ episodes $\times$ $40$ days.
\begin{itemize}
  \item \textbf{HARD (LRU eviction).} Heads occupy discrete slots ordered by
        recency of training. Training on a held regime refreshes its recency.
        Training on an \emph{unheld} regime at full capacity instantiates a
        new head and \textbf{evicts} the least-recently-used head
        \emph{together with its episode buffer}---competence and memory are
        destroyed jointly and completely.
  \item \textbf{SOFT (interference decay).} No head is ever evicted. Instead,
        every foreign update multiplies \emph{all other held heads} by
        $(1 - \rho_{\mathrm{int}})$ with $\rho_{\mathrm{int}} = 0.06$:
        after $N$ foreign updates a dormant head retains strength
        $(1 - \rho_{\mathrm{int}})^{N}$ of its trained weights.
\end{itemize}
\end{definition}

The hard model is the cache-eviction picture; the soft model is the
gradient-interference picture of catastrophic forgetting in neural networks
\citep{french1999catastrophic, kirkpatrick2017overcoming, parisi2019continual},
where training on task $B$ perturbs the very weights encoding task $A$.
Progressive networks \citep{rusu2016progressive} avoid interference by
freezing columns---a structural cousin of what pinning will achieve in
Section~\ref{sec:dd2-competition}.

\subsection{Hard Model: The Coupon-Collector Clock}

Under HARD, a dormant head dies when enough \emph{distinct} foreign regimes
have been learned to push it out of the LRU order. The clean question: a
learner holds the dormant head and has $K$ slots; how long until it has
trained on $K$ distinct foreign regimes (each claiming a slot, the $K$-th
overflowing capacity and evicting the dormant head, which---having received no
training all dormancy---is always the least recently used)?

\begin{proposition}[Expected time to eviction]
\label{prop:dd2-coupon}
Suppose each dormancy day is independently a foreign-update day with
probability $p$, and on an update day the active regime is uniform on the
$W_a$ foreign regimes, independently across update days. Let $T_K$ be the
number of days until $K$ distinct foreign regimes have been trained on
($K \le W_a$). Then
\begin{equation}
\label{eq:dd2-coupon}
\E[T_K] \;=\; \frac{1}{p} \sum_{j=0}^{K-1} \frac{W_a}{W_a - j}.
\end{equation}
\end{proposition}

\begin{proof}
Stage $j$ begins when exactly $j$ distinct foreign regimes have been
collected, $0 \le j \le K-1$. On each day of stage $j$, the event ``collect a
new regime today'' requires an update day (probability $p$) on which the
active regime is one of the $W_a - j$ uncollected ones (probability
$(W_a - j)/W_a$, by uniformity and independence). The stage-$j$ waiting time
is therefore geometric with success probability $p (W_a - j)/W_a$ and
expectation $W_a / \bigl(p (W_a - j)\bigr)$. The stages are traversed in
order, so by linearity of expectation
$\E[T_K] = \sum_{j=0}^{K-1} W_a / \bigl(p(W_a - j)\bigr)$, which is
\eqref{eq:dd2-coupon}.
\end{proof}

\begin{example}[Worked numbers: $W_a = 3$, $K = 2$, $p = \tfrac12$]
\label{ex:dd2-coupon}
A learner holds the dormant \textsf{crisis} head with capacity $K = 2$; the
$W_a = 3$ common regimes are active; it trains on the active regime every
other day on average. By \eqref{eq:dd2-coupon},
\[
\E[T_2]
= \frac{1}{0.5}\left(\frac{3}{3} + \frac{3}{2}\right)
= 2 \times 2.5
= 5 \text{ days}:
\]
$1/p = 2$ expected days for any foreign update to claim the spare slot, plus
$3/(0.5 \times 2) = 3$ for an update from one of the remaining two regimes.
Five days---against a crisis dormancy of $250$--$450$ days: the dormant
head's expected lifetime is roughly $1\%$ of the dormancy it must survive.
\end{example}

How certain is eviction by day $D$? For the worked example we can compute the
survival probability exactly.

\begin{proposition}[Survival probability, exact and bounded]
\label{prop:dd2-survival}
In the setting of Proposition~\ref{prop:dd2-coupon} with $W_a = 3$, $K = 2$:
\begin{equation}
\label{eq:dd2-survival}
\Pr(\text{head survives } D \text{ days})
= 3\Bigl(1 - \tfrac{2p}{3}\Bigr)^{D} - 2(1 - p)^{D}.
\end{equation}
In general ($K \le W_a$), a union bound over which $K-1$ regimes were the
only ones updated gives the binomial-tail bound
\[
\Pr(\text{survive } D)
\;\le\; \binom{W_a}{K-1}
\Bigl(1 - p\,\tfrac{W_a - K + 1}{W_a}\Bigr)^{D}.
\]
\end{proposition}

\begin{proof}
Survival through day $D$ means at most $K - 1 = 1$ distinct regime was
trained on. Condition on $N_D \sim \mathrm{Bin}(D, p)$ update days: given
$N_D = m \ge 1$, all $m$ uniform draws land on a single regime with
probability $3 \cdot (1/3)^m$; given $m = 0$, survival is certain. Hence
\[
\Pr(\text{survive})
= (1-p)^D + \sum_{m \ge 1} \binom{D}{m} p^m (1-p)^{D-m} \cdot 3 \cdot 3^{-m}
= (1-p)^D + 3\Bigl[\bigl(1 - p + \tfrac{p}{3}\bigr)^{D} - (1-p)^D\Bigr],
\]
which is \eqref{eq:dd2-survival}. General bound: survival requires the
updated regimes to lie in some $(K-1)$-subset $S$; for fixed $S$, each day
avoids ``update outside $S$'' with probability $1 - p(W_a - K + 1)/W_a$;
union over the $\binom{W_a}{K-1}$ subsets (this matches the leading term of
\eqref{eq:dd2-survival} at $W_a = 3$, $K = 2$).
\end{proof}

\begin{example}[Eviction probabilities by dormancy length]
With $p = 0.5$, evaluating \eqref{eq:dd2-survival} (note
$1 - 2p/3 = 2/3$, $1 - p = 1/2$):

\begin{center}
\begin{tabular}{rccc}
\toprule
$D$ (days) & $3(2/3)^D$ & $2(1/2)^D$ & $\Pr(\text{evicted by } D)$ \\
\midrule
5  & $0.39506$ & $0.06250$  & $1 - 0.33256 = 0.667$ \\
10 & $0.05202$ & $0.00195$  & $1 - 0.05007 = 0.950$ \\
20 & $0.00090$ & $1.9\times 10^{-6}$ & $1 - 0.00090 = 0.9991$ \\
40 & $2.7\times 10^{-7}$ & $1.8\times 10^{-12}$ & $1 - 2.7\times 10^{-7}$ \\
\bottomrule
\end{tabular}
\end{center}

By $D = 20$ days---half of one common-regime block---eviction is a $99.9\%$
event; at a crisis dormancy of $D = 250$ survival is
$3(2/3)^{250} \approx 3 \times 10^{-44}$. Under the hard model an
unprotected crisis head \emph{does not survive dormancy}, full stop. The
i.i.d.-day idealization is conservative here: the real stream arrives in
persistent blocks, each delivering many updates from one regime, so
distinct-regime collection is at least as fast at the block scale.
\end{example}

\begin{example}[LRU mechanics, step by step]
\label{ex:dd2-lru}
A capacity-$K{=}2$ learner ends a crisis episode holding heads for
$\{\textsf{crisis}, \textsf{trend}\}$, with \textsf{crisis} most recent.
The dormancy stream then activates \textsf{trend}, \textsf{range},
\textsf{highvol}, \ldots\ and the learner trains on whatever is active. Slot
contents (ordered LRU $\to$ MRU; each head carries its episode buffer):

\begin{center}
\begin{tabular}{clll}
\toprule
Event & Active regime & Action & Slots after (LRU $\to$ MRU) \\
\midrule
1 & \textsf{trend}   & hit: refresh recency & $[\textsf{crisis}, \textsf{trend}]$ \\
2 & \textsf{range}   & miss at full capacity: \textbf{evict \textsf{crisis}} + buffer & $[\textsf{trend}, \textsf{range}]$ \\
3 & \textsf{crisis} reactivates & miss: \textbf{cold start}, evict \textsf{trend} & $[\textsf{range}, \textsf{crisis}^{\,\text{new}}]$ \\
\bottomrule
\end{tabular}
\end{center}

At event 2---the \emph{first} foreign block from an unheld regime---the
crisis head and its entire episode library are gone. At event 3 the probe
window is served by a freshly initialized head: full cold-start regret, then
a 15--30-day scramble to relearn what was once known, and the cycle repeats.
\end{example}

\subsection{Soft Model: Geometric Erosion and the Half-Life}

Under SOFT nothing is evicted; instead each foreign update multiplies every
other held head by $(1 - \rho_{\mathrm{int}})$ with
$\rho_{\mathrm{int}} = 0.06$. After $N_D$ foreign updates during a dormancy,
the dormant head retains
\begin{equation}
\label{eq:dd2-soft-retention}
\text{retained strength} \;=\; (1 - \rho_{\mathrm{int}})^{N_D} \;=\; 0.94^{N_D},
\end{equation}
of its trained weights (the head decays toward initialization; it is never
refreshed because its regime emits no training days). The natural summary is
the \textbf{half-life in foreign updates}: the $N$ solving
$(1-\rho)^{N} = \tfrac12$,
\begin{equation}
\label{eq:dd2-halflife}
N_{1/2} \;=\; \frac{\log(1/2)}{\log(1 - \rho_{\mathrm{int}})}
\;=\; \frac{-0.693147}{\log(0.94)}
\;=\; \frac{0.693147}{0.061875}
\;\approx\; 11.2 \text{ foreign updates},
\end{equation}
using $\log(0.94) = -0.061875$. Eleven foreign updates---about two competition
windows' worth of foreign training, or a week and a half of daily monolith
training---halve the head. The full erosion schedule:

\begin{center}
\begin{tabular}{rcccccc}
\toprule
Foreign updates $N$ & 11 & 12 & 25 & 50 & 100 & 250 \\
\midrule
Retained $0.94^{N}$ & $0.506$ & $0.476$ & $0.213$ & $0.045$ & $0.0021$ & $1.9 \times 10^{-7}$ \\
\bottomrule
\end{tabular}
\end{center}

For a monolith that trains every day ($N_D = D$), a crisis dormancy of
$250$--$450$ days leaves $0.94^{250} \approx 2 \times 10^{-7}$ of the head:
a cold start without a discrete eviction event. A learner suffering one
foreign update per 20-day window over the same dormancy takes
$N_D \approx 12$--$22$ updates, retention $0.48$--$0.25$: degraded, not
destroyed. The soft model grades the damage by the \emph{training cadence}
of the allocation rule---exactly the quantity the eleven arms vary.

\begin{keypoint}[title={Two models, one phenomenon, bounded from both ends}]
HARD and SOFT bracket the same constraint---finite capacity under foreign
training---from opposite extremes. HARD is the worst case per event: one
overflow update destroys a head completely, at a stopping time whose law
Propositions~\ref{prop:dd2-coupon}--\ref{prop:dd2-survival} pin down. SOFT is
the gentlest plausible case: no single update matters, yet the geometric
decay \eqref{eq:dd2-soft-retention} reaches the same endpoint at the
dormancies in play ($0.94^{250} \approx 10^{-7}$). A conclusion that holds
under \emph{both} is about bounded capacity itself, not a bookkeeping rule;
and a mechanism that protects a head must both prevent overflow updates
(HARD) and drive the foreign-update count to zero (SOFT). Pinning
(Section~\ref{sec:dd2-competition}) does both at once: after a covering
one-per-regime assignment, each specialist's foreign-update probability is
structurally $p = 0$.
\end{keypoint}

\begin{remark}[Why $\Phi$ is defined against the retention counterfactual]
The forgetting deficit \eqref{eq:dd2-phi} subtracts the regret of the
arm's \emph{own frozen past self}, not of an oracle, isolating dormancy
damage from competence: a mediocre head retained perfectly has $\Phi = 0$
(a skill problem, not a memory problem); an excellent head lost to eviction
has large $\Phi$ even at moderate probe regret. The eleven-arm comparison of
Section~\ref{sec:dd2-arms} needs both coordinates---probe regret for the
bottom line, $\Phi$ for the diagnosis.
\end{remark}

% ============================================================================
% PART III: THE RISP MECHANISM
% ============================================================================
\part{The \nm{} Mechanism}
\label{sec:dd2-part-mechanism}

Part~II established the problem: regimes recur after dormancies that no
capacity-limited, reward-chasing learner survives, while episode
heterogeneity punishes whoever survives by memorizing too literally. This
part builds the mechanism that answers it: specialists
(Section~\ref{sec:dd2-specialist}), windowed competition with multiplicative
affinity updates and an absorbing pin (Section~\ref{sec:dd2-competition}),
and the placement of \nm{} among ten alternatives in a two-axis design space
(Section~\ref{sec:dd2-arms}).

\section{The Specialist}
\label{sec:dd2-specialist}

\subsection{State: Identity, Competence, Memory}

\begin{definition}[Specialist state]
\label{def:dd2-specialist}
The population contains $N = 4$ specialists. Specialist $i$ at day $t$
carries three coupled state components:
\begin{enumerate}
  \item \textbf{Identity: niche affinity} $\alpha_i \in \Delta^{R}$,
        initialized uniform ($\alpha_{i,r} = 1/R = 0.25$): the specialist's
        \emph{claim} on each regime, updated only by winning competition
        windows (Section~\ref{sec:dd2-competition}); it controls the
        incumbency bonus and, ultimately, pinning.
  \item \textbf{Competence: regime heads} $\{w_{i,r} : r \in H_i(t)\}$, one
        linear SPO+ head per \emph{held} regime, $|H_i(t)| \le K < R$,
        mapping features to predicted returns
        $\hat y_{t,i'} = w_{i,r} \cdot x_{t,i'}$ and thence to a decision
        via the optimizer (Definition~\ref{def:dd2-regret}).
  \item \textbf{Memory: episode library}
        $\{\mathcal{B}_{i,r} : r \in H_i(t)\}$, at most $6$ episodes
        $\times$ $40$ days per held regime of $(X_t, y_t)$ records keyed by
        episode---created with the head and \textbf{evicted with the head}
        (hard model); the raw material for both training variants below.
\end{enumerate}
\end{definition}

The three components are deliberately welded together: affinity determines
where the specialist wins; winning triggers training, which builds the head
and writes the buffer; head and buffer determine future window performance,
which feeds back into affinity. Identity, competence, and memory rise
together and (under eviction) fall together.

\begin{definition}[Decision regret]
\label{def:dd2-regret}
Fix a compact feasible decision set $\mathcal{A} \subset \R^{n}$ (a polytope
of admissible portfolios over the $n = 20$ assets), and let
$a^{*}(v) \in \argmax_{a \in \mathcal{A}} v^{\top} a$ with optimal value
$z^{*}(v) = \max_{a \in \mathcal{A}} v^{\top} a$. The \textbf{per-day
decision regret} of a prediction $\hat y_t$ against realized returns $y_t$ is
\begin{equation}
\label{eq:dd2-regret}
\rho(\hat y_t,\, y_t) \;=\; z^{*}(y_t) \;-\; y_t^{\top} a^{*}(\hat y_t)
\;\ge\; 0:
\end{equation}
the return foregone by acting on the prediction instead of on the truth.
Regret depends on $\hat y_t$ only through the \emph{decision} it induces---two
predictions with different numbers but the same induced portfolio have
identical regret. This is the decision-focused viewpoint of
\citet{elmachtoub2022smart}: the loss is downstream decision quality, not
prediction error.
\end{definition}

\subsection{One Day in the Life}

Each day $t$, the causal labeler emits $\hat r_t$ (visible to everyone), and
the system follows one of two paths.

\textbf{Pinned case} ($\hat r_t$ pinned to owner $o(\hat r_t)$,
Definition~\ref{def:dd2-pin}):
\begin{enumerate}
  \item The owner alone serves the live decision
        $a_t = a^{*}\bigl(\hat y^{(o)}_t\bigr)$ using $w_{o, \hat r_t}$.
  \item The day's record $(X_t, y_t)$ is appended to the owner's buffer
        $\mathcal{B}_{o, \hat r_t}$ (current episode, capped at $40$ days;
        oldest of the $\le 6$ episodes rotated out when a new one opens).
  \item On its training cadence, the owner takes a training step
        (Section~\ref{sec:dd2-training}) on $\mathcal{B}_{o, \hat r_t}$.
        \emph{No other specialist decides, trains, or updates anything.}
\end{enumerate}

\textbf{Unpinned case}:
\begin{enumerate}
  \item \textbf{Shadow decisions.} Every specialist $i$ forms
        $\hat y^{(i)}_t$ (with $w_{i, \hat r_t}$ if held; with the
        zero-initialized default head, hence a default decision, if not) and
        its shadow decision $a^{(i)}_t = a^{*}(\hat y^{(i)}_t)$. The live
        decision is served by the current best claimant (highest
        $\alpha_{\cdot, \hat r_t}$).
  \item \textbf{Regret accounting.} Each specialist's realized regret
        $\rho^{(i)}_t = \rho(\hat y^{(i)}_t, y_t)$ is appended to the open
        competition window for regime $\hat r_t$. Windows are regime-gated:
        only days with $\hat r_t = r$ advance $r$'s window, so a window
        spans regime interruptions without contamination.
  \item \textbf{Window close.} When $r$'s window reaches $W_c = 20$
        in-regime days, a winner is declared
        (Section~\ref{sec:dd2-competition}), the winner's affinity takes one
        exponentiated-gradient step, and \emph{the winner alone} trains on
        the window's data (appended to its buffer for $r$; if it holds no
        head for $r$, one is instantiated now---the moment HARD eviction or
        SOFT interference strikes its \emph{other} heads).
  \item \textbf{Pin check.} If $\alpha_{i^{*}, r} \ge 0.95$, regime $r$ is
        pinned to $i^{*}$ permanently.
\end{enumerate}

\begin{intuition}[title={Why losers do not train}]
The most consequential line in the loop is the quiet one: only the window
winner trains. Training is simultaneously the reward (the winner improves at
$r$) and the cost-bearer (training on $r$ is a foreign update against every
other head the trainer holds). A specialist that keeps losing on $r$ never
pays $r$'s interference cost---its own niche's head sits at $p = 0$ in the
foreign-update process of Section~\ref{sec:dd2-dormancy}. Competition is not
just selecting who improves; it is rationing who gets \emph{damaged}.
\end{intuition}

\subsection{The SPO+ Training Step}
\label{sec:dd2-training}

Training must improve \emph{decisions}, so the loss is a decision-aware
surrogate: regret \eqref{eq:dd2-regret} itself is piecewise constant in
$\hat y$ (it changes only when the induced $\argmax$ jumps), hence useless
for gradients. The SPO+ surrogate of \citet{elmachtoub2022smart}, written
here in the return-maximization convention, is
\begin{equation}
\label{eq:dd2-spoloss}
\ell_{\mathrm{SPO+}}(\hat y, y)
\;=\;
\max_{a \in \mathcal{A}} \bigl\{ (2\hat y - y)^{\top} a \bigr\}
\;-\; 2 \hat y^{\top} a^{*}(y)
\;+\; z^{*}(y).
\end{equation}

\begin{lemma}[SPO+ is a convex upper bound on regret]
\label{lem:dd2-spoplus}
$\ell_{\mathrm{SPO+}}(\cdot, y)$ is convex in $\hat y$, satisfies
$\ell_{\mathrm{SPO+}}(y, y) = 0$, and
$\ell_{\mathrm{SPO+}}(\hat y, y) \ge \rho(\hat y, y)$ for all $\hat y$.
A subgradient in $\hat y$ is
\begin{equation}
\label{eq:dd2-spograd}
g(\hat y, y) \;=\; 2\bigl(a^{*}(2\hat y - y) - a^{*}(y)\bigr).
\end{equation}
\end{lemma}

\begin{proof}
Convexity: the first term of \eqref{eq:dd2-spoloss} is a maximum of affine
functions of $\hat y$, the rest linear or constant. At $\hat y = y$ the
three terms are $z^{*}(y) - 2z^{*}(y) + z^{*}(y) = 0$. Upper bound:
evaluating the max at the feasible point $a^{*}(\hat y)$,
\[
\ell_{\mathrm{SPO+}}(\hat y, y)
\;\ge\; 2\hat y^{\top}\bigl(a^{*}(\hat y) - a^{*}(y)\bigr)
+ \bigl(z^{*}(y) - y^{\top} a^{*}(\hat y)\bigr),
\]
where the first bracket is $\ge 0$ because $a^{*}(\hat y)$ maximizes
$\hat y^{\top} a$, and the second is exactly $\rho(\hat y, y)$.
Subgradient: Danskin's theorem on the max term gives
$2\,a^{*}(2\hat y - y)$; the linear term contributes $-2a^{*}(y)$.
\end{proof}

With the linear head $\hat y_{t,i} = w^{\top} x_{t,i}$, the chain rule turns
\eqref{eq:dd2-spograd} into a gradient in $w$:
$\nabla_{w}\, \ell_{\mathrm{SPO+}}(\hat y_t, y_t)
= \sum_{i=1}^{n} g_i(\hat y_t, y_t)\, x_{t,i}$.

\begin{example}[SPO+ arithmetic on two assets]
\label{ex:dd2-spoplus}
Let $\mathcal{A} = \{e_1, e_2\}$, truth $y = (0.01, -0.01)$, prediction
$\hat y = (-0.002, 0.004)$ (ranking inverted). Then $a^{*}(y) = e_1$,
$z^{*}(y) = 0.01$, $a^{*}(\hat y) = e_2$, so
$\rho = 0.01 - (-0.01) = 0.02$. For SPO+:
$2\hat y - y = (-0.014, 0.018)$, maximized by $e_2$ at $0.018$; the linear
term is $-2\hat y^{\top} e_1 = 0.004$; hence
$\ell_{\mathrm{SPO+}} = 0.018 + 0.004 + 0.01 = 0.032 \ge \rho = 0.02$.
The subgradient \eqref{eq:dd2-spograd} is $g = 2(e_2 - e_1) = (-2, 2)$:
descent pushes $\hat y_1$ up and $\hat y_2$ down---directly toward repairing
the ranking that caused the bad decision, not toward shrinking prediction
error per se.
\end{example}

\subsection{ERM versus INV over the Episode Library}
\label{sec:dd2-erm-inv}

A training step consumes the episode buffer
$\mathcal{B}_{i,r} = \{B_1, \ldots, B_E\}$, $E \le 6$ episodes. Write
\[
\bar\ell_e(w) \;=\; \frac{1}{|B_e|} \sum_{t \in B_e}
\ell_{\mathrm{SPO+}}\bigl(\hat y_t(w),\, y_t\bigr),
\qquad
\bar\ell(w) \;=\; \frac{1}{E} \sum_{e=1}^{E} \bar\ell_e(w)
\]
for the per-episode and pooled mean surrogate losses. The two variants:

\paragraph{ERM variant.} Sample one episode $e \sim \mathrm{Unif}\{1..E\}$
and take a subgradient step on $\bar\ell_e(w)$. In expectation this descends
the pooled objective $\bar\ell(w)$: all episode-days are equal, and the head
converges toward the buffer-pooled SPO+ minimizer---including whatever
episode-specific $\gamma$ structure dominates the buffer.

\paragraph{INV variant.} Descend the variance-penalized objective
\begin{equation}
\label{eq:dd2-invobj}
J(w) \;=\; \Var_e\bigl[\bar\ell_e(w)\bigr] \;+\; \beta\, \E_e\bigl[\bar\ell_e(w)\bigr]
\;=\; \frac{1}{E} \sum_{e=1}^{E} \bigl(\bar\ell_e(w) - \bar\ell(w)\bigr)^2
\;+\; \beta\, \bar\ell(w),
\end{equation}
which demands not just low mean loss but \emph{evenly} low loss across
episodes---the invariance principle transplanted into decision-focused
learning \citep{yan2026invpnco}.

\begin{proposition}[INV gradient]
\label{prop:dd2-invgrad}
The gradient of \eqref{eq:dd2-invobj} is
\begin{equation}
\label{eq:dd2-invgrad}
\nabla_w J(w) \;=\; \frac{1}{E} \sum_{e=1}^{E}
\Bigl[\, 2\bigl(\bar\ell_e(w) - \bar\ell(w)\bigr) + \beta \,\Bigr]\,
\nabla_w \bar\ell_e(w).
\end{equation}
\end{proposition}

\begin{proof}
Differentiate the variance term of \eqref{eq:dd2-invobj}:
\[
\nabla_w \Var_e
= \frac{1}{E} \sum_{e} 2\bigl(\bar\ell_e - \bar\ell\bigr)
\bigl(\nabla \bar\ell_e - \nabla \bar\ell\bigr)
= \frac{2}{E} \sum_{e} \bigl(\bar\ell_e - \bar\ell\bigr) \nabla \bar\ell_e
\;-\; \frac{2}{E} \nabla \bar\ell \underbrace{\sum_{e} \bigl(\bar\ell_e - \bar\ell\bigr)}_{=\,0}.
\]
The cross term vanishes identically because the centered residuals
$\bar\ell_e - \bar\ell$ sum to zero by definition of $\bar\ell$---this is the
one step worth seeing explicitly, as it is why the INV gradient is a clean
\emph{reweighting} of per-episode gradients rather than a new object. Adding
$\beta \nabla \bar\ell = (\beta/E) \sum_e \nabla \bar\ell_e$ and collecting
terms gives \eqref{eq:dd2-invgrad}.
\end{proof}

\begin{intuition}[title={Reading the INV weights}]
Equation~\eqref{eq:dd2-invgrad} gives episode $e$ the gradient weight
$2(\bar\ell_e - \bar\ell) + \beta$: an episode at the buffer average gets
weight $\beta$ (plain ERM); a worse-than-average episode gets weight
$> \beta$; and a better-than-average episode with
$\bar\ell_e < \bar\ell - \beta/2$ gets a \emph{negative} weight---gradient
\emph{ascent} on episodes served suspiciously well. That last case is the
anti-overfitting tooth. In the generator \eqref{eq:dd2-generator}, loading
on the stable direction $\theta_r$ lowers all $\bar\ell_e$ roughly equally
(no variance contribution), while loading on $x_{\mathrm{sp}}$ helps the
episodes whose $\gamma_{r,e}$ sign matches and hurts the rest---pumping the
variance term. INV strips $\gamma$-chasing weight while keeping
$\theta$-weight: the objective-level answer to the het trap of
Example~\ref{ex:dd2-het}, complementary to the allocation-level answer
(pinning). At $\mathrm{het} = 0$ the penalty has nothing real to penalize
and INV pays a small price---which is why the experiment carries both A5
(ERM) and A6 (INV) rather than declaring a winner a priori.
\end{intuition}

\section{Competition Dynamics}
\label{sec:dd2-competition}

\subsection{Windows, Winners, and the Exponentiated-Gradient Update}

\begin{definition}[Competition window and winner]
\label{def:dd2-window}
For an unpinned regime $r$, the competition window collects, for each
specialist $i$, its shadow decision regrets over $W_c = 20$ in-regime days.
At window close the \textbf{winner} is
\begin{equation}
\label{eq:dd2-winner}
i^{*} \;=\; \argmin_{i \in \{1..N\}} \;
\Bigl[\, \bar\rho^{(i)}_{\mathrm{win}}
\;-\; \lambda \bigl(\alpha_{i,r} - \tfrac{1}{R}\bigr) \Bigr],
\qquad \lambda = 0.3,
\end{equation}
where $\bar\rho^{(i)}_{\mathrm{win}}$ is $i$'s mean window regret: lowest
regret wins, with an \textbf{incumbency bonus} $\lambda(\alpha_{i,r} - 1/R)$
subtracted for affinity above uniform. The winner then applies one
\textbf{exponentiated-gradient (EG) step} to its own affinity vector,
\begin{equation}
\label{eq:dd2-eg}
\alpha_{i^{*}, r} \;\leftarrow\; \alpha_{i^{*}, r}\, e^{\eta},
\qquad \text{then renormalize } \alpha_{i^{*}} \text{ over } \Delta^{R},
\qquad \eta = 0.6,
\end{equation}
and trains on the window's data. Losers' affinity vectors are untouched.
\end{definition}

Two structural remarks. First, the bonus is centered at $1/R$, so an
uncommitted specialist gets zero bonus, and since
$\sum_r (\alpha_{i,r} - 1/R) = 0$, incumbency advantages across regimes are
\emph{zero-sum}: claiming $r$ harder structurally weakens the claim
everywhere else. Second, \eqref{eq:dd2-eg} belongs to the
multiplicative-update family of exponentiated gradient
\citep{kivinen1997exponentiated}, Hedge \citep{freund1997decision}, and
multiplicative weights \citep{arora2012multiplicative}---inverted: each
\emph{specialist} runs the update over \emph{regimes} (what shall I be?),
rather than a master algorithm weighting experts (whom shall I trust?).

Because only the $r$-coordinate is multiplied and the rest are renormalized
proportionally, the update has a two-point reduction: writing
$\alpha = \alpha_{i^{*}, r}$ before the win,
\begin{equation}
\label{eq:dd2-egclosed}
\alpha' \;=\; \frac{\alpha\, e^{\eta}}{\alpha\, e^{\eta} + (1 - \alpha)},
\qquad \text{equivalently} \qquad
\underbrace{\log \frac{\alpha'}{1 - \alpha'}}_{\text{log-odds}}
= \log \frac{\alpha}{1 - \alpha} + \eta .
\end{equation}
Each win adds exactly $\eta$ to the log-odds of regime $r$ within the
winner's identity. Affinity is a win counter, read through a sigmoid.

\begin{example}[Worked EG trace: uniform to pinned in seven wins]
\label{ex:dd2-egtrace}
Start at $\alpha = 1/R = 0.25$ and let the same specialist win consecutive
windows for $r$ (winning nothing else in between). With $\eta = 0.6$,
$e^{0.6} = 1.82212$. First win, by \eqref{eq:dd2-egclosed}:
\[
\alpha_1 = \frac{0.25 \times 1.82212}{0.25 \times 1.82212 + 0.75}
= \frac{0.45553}{1.20553} = 0.37787.
\]
Second win:
\[
\alpha_2 = \frac{0.37787 \times 1.82212}{0.37787 \times 1.82212 + 0.62213}
= \frac{0.68852}{1.31065} = 0.52533.
\]
Iterating (all values rounded from full-precision iteration):

\begin{center}
\begin{tabular}{lccccccccc}
\toprule
Wins $k$ & 0 & 1 & 2 & 3 & 4 & 5 & 6 & 7 & 8 \\
\midrule
$\alpha_k$ & $0.250$ & $0.378$ & $0.525$ & $0.668$ & $0.786$ & $0.870$ & $0.924$ & $\mathbf{0.957}$ & $0.976$ \\
\bottomrule
\end{tabular}
\end{center}

The pin threshold $0.95$ is crossed at win $7$. The closed form confirms it:
$\alpha_k \ge 0.95$ iff the odds satisfy
$\tfrac{\alpha_0}{1-\alpha_0} e^{\eta k} = \tfrac13 e^{0.6k} \ge 19$, i.e.\
$e^{0.6k} \ge 57$, i.e.\ $k \ge \ln 57 / 0.6 = 4.043/0.6 = 6.74$, so
$k = 7$. Seven clean wins---$7 \times 20 = 140$ in-regime days, roughly three
to four common-regime episodes---separate a cold uniform specialist from
permanent ownership.
\end{example}

\subsection{Rich-Get-Richer: The Symmetry-Breaking Engine}

The EG update alone would make ownership a random walk among near-identical
specialists. What breaks the symmetry is the coupling to training:
\emph{winners train $\Rightarrow$ training improves the head and deepens the
buffer $\Rightarrow$ a better head wins the next window more often
$\Rightarrow$ winners win again.}

Formally, let $q_m$ be the probability that the current leader on regime $r$
wins window $m$. Early on, heads are near-identical and $q_m \approx$ the
tie-breaking rate, but the first (possibly lucky) win hands the winner a
window's worth of training on $r$; under the separated-$\theta_r$ generator
this lowers its expected window regret, so $q_{m+1} > q_m$; and the affinity
bonus in \eqref{eq:dd2-winner} adds a deterministic tilt
$\lambda(\alpha_{i,r} - \alpha_{j,r})$ on top. The win process is a
reinforcing urn: the leader's log-odds climb at rate $\eta$ per win, while
losers---who do not train on $r$---suffer no foreign updates from it. Once
one specialist leads on each regime, the system is in the basin of the
covering one-per-regime assignment.

This is competitive exclusion in the sense of \citet{li2026gause}: one
species per niche at equilibrium, with reward competition as the exclusion
force---and the empirical record matches the mechanism's prediction exactly:
in \textbf{all 20 seeds}, the population converges to a covering
one-specialist-per-regime assignment, with each common regime pinned at
$\alpha \ge 0.95$.

\subsection{The Batching Question: A Pre-Registered Prediction and Its Refutation}
\label{sec:dd2-batching}

How large must $W_c$ be? The pre-registered reasoning was a single-window
concentration argument. Suppose two specialists' true mean per-day regrets
on $r$ differ by $\Delta > 0$, with per-day regrets bounded in $[0, B]$. The
window comparison errs when the worse specialist posts the lower sample
mean; by Hoeffding's inequality on the per-day regret differences (range
$2B$, mean $\Delta$),
\begin{equation}
\label{eq:dd2-hoeffding}
\Pr\bigl(\text{wrong winner}\bigr)
\;\le\; \exp\!\left( - \frac{W_c\, \Delta^2}{2 B^2} \right),
\qquad \text{so} \qquad
W_c \;\ge\; \frac{2 B^2}{\Delta^2} \ln \frac{1}{\delta}
\end{equation}
guarantees error $\le \delta$. Early in competition $\Delta$ is small by
construction (symmetric initialization), so \eqref{eq:dd2-hoeffding} blows
up, and the natural prediction follows: \emph{small windows produce noisy
winners, noisy winners scatter the EG updates, and specialization should
require large $W_c$}---batching looked necessary, with $W_c = 1$ expected to
fail outright.

\begin{keypoint}[title={The empirical refutation, stated plainly}]
The pre-registered prediction was wrong. Sweeping
$W_c \in \{1, 5, 20, 60\}$ over 20 seeds, post-reactivation regret is
\textbf{flat}: $0.0303$--$0.0309$ across the sweep, with $W_c = 1$
\emph{marginally best}. Batching within the window is not necessary for
convergence, and the analysis that said otherwise mislocated where the
averaging happens.
\end{keypoint}

The corrected mechanism story. Inequality \eqref{eq:dd2-hoeffding} bounds the
error of a \emph{single} comparison, but nothing in the architecture consumes
single comparisons: the EG recursion \eqref{eq:dd2-egclosed} sums
$\eta \cdot \mathbf{1}\{\text{win}\}$ in log-odds space across windows, and
reaching the pin requires $\approx \ln 57 / \eta \approx 6.7$ \emph{net}
wins, not one correct vote. If the better specialist wins each window with
probability $q$ above its rivals' rates, its lead in win counts---hence in
log-odds---grows linearly in the number of windows with noise
$O(\sqrt{\text{windows}})$: the law of large numbers operates \emph{across}
windows, in $\alpha$-space, regardless of $W_c$. Shrinking $W_c$ from 20 to 1
makes each vote noisier but delivers $20\times$ more votes per regime-day,
and the two effects approximately cancel; small windows also start the
rich-get-richer training loop earlier (the first win arrives after 1 day, not
20), the plausible source of $W_c = 1$'s marginal edge. The affinity is
already an exponentially-weighted tally of the win history---the EG
accumulation \emph{is} the cross-window averaging---so a second averaging
layer inside the window was redundant.

\begin{warning}[title={Do not quietly upgrade the corrected story to a prediction}]
The flat $W_c$ sweep is an empirical result on this stream (synthetic
generator, $R = 4$, $N = 4$, $\eta = 0.6$), not a theorem. The corrected
account is a post-hoc narrative consistent with the data and with standard
multiplicative-weights analyses
\citep{freund1997decision, arora2012multiplicative}, but the experiment that
would test \emph{it} (varying $\eta$ jointly with $W_c$) was not
pre-registered. One box above records a refuted prediction; this box keeps
the replacement story from enjoying a credibility it has not yet earned.
\end{warning}

\subsection{Pinning}
\label{sec:dd2-pinning}

\begin{definition}[Pin]
\label{def:dd2-pin}
When a window winner's updated affinity satisfies
$\alpha_{i^{*}, r} \ge 0.95$, regime $r$ is \textbf{pinned} to owner
$o(r) = i^{*}$, permanently: thereafter $r$'s live decisions are served by
$o(r)$ alone, training on $r$-days goes to $o(r)$ alone, and the competition
machinery for $r$ (shadow decisions, windows, EG updates) shuts down.
\end{definition}

What changes at the pin is the \emph{modality} of protection. Before it, the
owner-apparent is protected statistically: it keeps winning, so it keeps
being the one who trains, but a bad run could in principle let a rival in.
After it, protection is structural: ownership no longer consults the reward
stream at all---the assignment becomes reward-independent, not as an
assumption but as a phase the dynamics enter and cannot leave. Once all
regimes are pinned one-per-specialist (the observed endpoint in every seed),
each specialist trains only on its own regime: $p = 0$, so under HARD it
never overflows and under SOFT it suffers zero interference. The
coupon-collector clock of Proposition~\ref{prop:dd2-coupon} never starts; the
half-life \eqref{eq:dd2-halflife} never ticks. Dormancy becomes silence, with
the head and the six-episode library untouched---a distributed persistent
memory.

Two ablations calibrate how much of the benefit is the pin itself versus the
competition that produces it:
\begin{itemize}
  \item \textbf{Never-pin:} disabling Definition~\ref{def:dd2-pin} costs
        $0.0320$ versus $0.0307$ pinned post-reactivation regret---the
        residual price of merely statistical protection: perpetual shadow
        competition occasionally hands windows, hence foreign training, to
        non-owners, and serving can flicker at reactivations.
  \item \textbf{No incumbency bonus:} $\lambda = 0$ yields $0.0314$ versus
        $0.0307$. The bonus \emph{accelerates} convergence to ownership but
        does not cause it---the rich-get-richer loop breaks symmetry on its
        own. Causality lives in the win--train coupling; $\lambda$ is a
        catalyst.
\end{itemize}

\paragraph{The crisis wrinkle.}
The rare regime exposes a quantization effect in the pin rule. A crisis
episode lasts $15$--$30$ days and yields at most one complete $W_c = 20$
window---often zero---and crisis episodes are themselves few
(Section~\ref{sec:dd2-scarcity}). Across the 20 seeds, the crisis regime's
top affinity reaches only $\alpha \in [0.67, 0.87]$: by the trace of
Example~\ref{ex:dd2-egtrace}, exactly $3$--$5$ EG wins from uniform
($\alpha_3 = 0.668$, $\alpha_5 = 0.870$)---the affinity arithmetic
independently confirms that the stream does not contain the $\approx 7$
complete crisis windows formal pinning requires. The substantive outcome is
nonetheless clean: in every seed the crisis regime has a \emph{stable unique
owner}, so the covering assignment is achieved; only the formal threshold is
starved. The deployment lesson: pin on \emph{ownership stability} (the same
specialist winning a regime's windows across $k$ consecutive episodes), not
on the raw $\alpha \ge 0.95$ threshold, whose win-count requirement assumes
window-rich regimes. For rare regimes, windows---not days, not
episodes---are the binding currency.

\section{The Eleven Arms as Points in Design Space}
\label{sec:dd2-arms}

The companion experiment runs eleven arms. They are not eleven competing
architectures; they are a factorial probe of a two-axis design space, and
each arm exists to occupy a cell or test an edge.

\paragraph{Axis 1: allocation reward-coupling.} How does the regime
$\to$ learner assignment respond to the \emph{current} reward stream?
Three positions: \textbf{coupled} (chases recent performance forever),
\textbf{transient-then-structural} (\nm{}: coupled during competition,
independent after the pin), and \textbf{independent} (lottery or oracle).
Part~II showed why this axis governs forgetting: coupled allocation keeps
every learner's foreign-update probability $p > 0$ during dormancy, and both
memory models convert $p > 0$ over $250$--$450$ days into a destroyed head.

\paragraph{Axis 2: training objective.} Pooled \textbf{ERM} or the
variance-penalized \textbf{INV} of \eqref{eq:dd2-invobj} on the buffered
episodes. This axis governs the het trap of Example~\ref{ex:dd2-het}: ERM
imports episode-specific $\gamma$-weight into the probe window; INV strips
it.

\begin{center}
\begin{tabular}{llll}
\toprule
Arm & Allocation rule & Coupling & Objective \\
\midrule
A1 & monolith: train on whatever is active & coupled (degenerate) & ERM \\
A2 & MoE learned router & coupled & ERM \\
A3 & recent-performance capital shares & coupled (smooth) & ERM \\
A4 & random fixed niches & independent & ERM \\
A5 & \nm{} competition $\to$ pin & transient $\to$ structural & ERM \\
A6 & \nm{} competition $\to$ pin & transient $\to$ structural & INV \\
A7 & monolith & coupled (degenerate) & INV \\
A8a & Hedge over fixed strategies & coupled (weights only) & none (no training) \\
A8b & Hedge over learning experts & coupled & ERM \\
A9 & oracle assignment from day 1 & independent & ERM \\
A10 & oracle assignment from day 1 & independent & INV \\
\bottomrule
\end{tabular}
\end{center}

\paragraph{A1 (monolith-ERM).} One capacity-$K$ learner; the allocation rule
is the degenerate \emph{always me}: it serves and trains on whatever
$\hat r_t$ is active. Its foreign-update process during any dormancy has
$p = 1$: under HARD it is the learner of Example~\ref{ex:dd2-lru}, evicted
within one foreign block; under SOFT it retains
$0.94^{250} \approx 2 \times 10^{-7}$ over a crisis dormancy.
\emph{Prediction:} catastrophic forgetting under both memory models;
worst-tier probe regret.

\paragraph{A2 (MoE learned router).} $N$ experts plus a gate matrix
$G \in \R^{R \times N}$; on each day the gate serves expert
$j_t = \argmax_j G[\hat r_t, j]$ with probability $0.9$ and a uniformly
random expert with probability $\epsilon = 0.1$ ($\epsilon$-greedy); the
served expert trains on the day; the gate updates
\[
G[\hat r_t, j_t] \;\mathrel{+}=\; \mathrm{lr} \cdot
\bigl(\text{baseline} - \rho_{t}\bigr),
\]
with the baseline a running mean regret---below-baseline performance raises
the routing weight. This is reward-coupled allocation \emph{forever}, and
exploration alone guarantees damage: during a crisis dormancy of $D$ days
the crisis expert is dragged into foreign training at rate
$\ge \epsilon/N$ per day, so
$\E[\text{foreign updates}] \ge \epsilon D / N
= 0.1 \times (250\text{--}450)/4 = 6.25\text{--}11.25$---about one full
half-life \eqref{eq:dd2-halflife} before counting routing churn; under HARD
a single explored foreign block can evict. MoE routing in continual settings
is studied in \citet{li2024moecl}; the failure here is not expressiveness
but the absence of any reward-independent protection phase.
\emph{Prediction:} forgets despite having one expert per regime available.

\paragraph{A3 (recent-performance allocator).} All learners persist; capital
shares are a softmax of recent performance, and each learner's
\emph{training intensity equals its share}. The instructive feature is the
\textbf{softmax floor}: a softmax is strictly positive, so every learner
receives a minimum share $s_{\min} > 0$ of every day's training, foreign
days included. The floor \emph{is} the erosion rate: a dormant regime's
natural owner accrues foreign updates at rate $\ge s_{\min}$ per day,
retention $\approx 0.94^{\,s_{\min} D}$ under SOFT---graded and unavoidable.
\emph{Prediction:} graded (not cliff-like) forgetting; better than A1/A2,
worse than any covering reward-independent arm.

\paragraph{A4 (random fixed niches).} Each specialist is assigned a regime
by lottery at $t = 0$, fixed forever: perfectly reward-independent, zero
foreign updates, $\Phi \approx 0$ on every \emph{covered} regime. The defect
is coverage: with $N = R = 4$ and independent uniform draws, all four
regimes are covered with probability $4!/4^4 = 24/256 \approx 9.4\%$---in
over $90\%$ of lotteries some regime has no owner and is served at
default-decision quality forever. Reward-independence is necessary but not
sufficient: the assignment must also cover.
\emph{Prediction:} bimodal---oracle-grade retention on covered regimes,
worst-case service on uncovered ones.

\paragraph{A5 / A6 (\nm{}-ERM / \nm{}-INV).} The mechanism of
Sections~\ref{sec:dd2-specialist}--\ref{sec:dd2-competition}, differing only
in the buffer objective (ERM step vs.\ INV gradient \eqref{eq:dd2-invgrad}).
\nm{} sits at the distinguished point of axis 1: reward-coupled
\emph{exactly long enough} to discover a covering assignment (solving A4's
coverage failure), then structurally independent (solving A1--A3's
protection failure). \emph{Predictions:} both converge to covering pinned
assignments and retain through dormancy (confirmed: 20/20 seeds, common
regimes pinned $\ge 0.95$, pinned regret $0.0307$); A6 outperforms A5 on
probe windows as $\mathrm{het}$ grows, since the probe meets a fresh
$\gamma_{r,e}$ (Example~\ref{ex:dd2-het}).

\paragraph{A7 (monolith-INV).} A1's allocation with A6's objective, blocking
a confound: if \nm{}-INV beats the monolith, is it the mechanism or just the
better loss? INV reshapes \emph{what} is learned from the buffer, but the
buffer itself is evicted (HARD) or the head eroded (SOFT) by the allocation;
no objective can preserve parameters it does not control.
\emph{Prediction:} A7 $\approx$ A1 on retention---the objective axis cannot
rescue a failure on the allocation axis.

\paragraph{A8a / A8b (Hedge over fixed / learning experts).} A8a runs Hedge
\citep{freund1997decision} over a pool of \emph{fixed} strategies: weights
$w_i \propto \exp(-\eta_H \sum_s \rho^{(i)}_s)$, no training anywhere. The
allocation is reward-coupled, but with nothing trainable there is nothing to
forget: A8a is structurally immune to dormancy, and bounded by the best
fixed strategy rather than the best learned head. A8b keeps the Hedge
allocation but lets the weighted experts train on the days they are trusted
with: trainability returns, and with it the foreign-update process.
\emph{A8a vs.\ A8b isolates whether dormancy damage comes from
reward-coupled allocation per se or from training under it} (the latter).
\emph{Predictions:} A8a immune to forgetting but capped within episodes;
A8b recovers in-episode adaptation and re-imports forgetting.

\paragraph{A9 / A10 (oracle assignment, ERM / INV).} The regime $\to$
specialist map is given correctly at $t = 0$---pinning before day one---and
each owner trains only on its regime, with ERM (A9) or INV (A10). These are
the ceilings: reward-independent, covering, zero discovery cost. The
A5/A6 vs.\ A9/A10 gap prices the \emph{discovery phase} (contested serving,
occasional wrong winners, foreign training before ownership settles); the
A9 vs.\ A10 gap reads off the pure objective effect with allocation held
perfect, mirroring A5 vs.\ A6.

\begin{keypoint}[title={The arms are an experiment about classes, not architectures}]
Each arm represents a \emph{cell} in (coupling $\times$ objective) space,
and every headline conclusion is a contrast between cells with one axis held
fixed: A5 vs.\ A2/A3 tests allocation class at fixed objective; A6 vs.\ A5
and A10 vs.\ A9 test objective class at fixed allocation; A7 shows the
objective axis cannot compensate for the allocation axis; A8a vs.\ A8b shows
the damage mechanism is training-under-coupling, not coupling itself; A4
shows reward-independence without coverage is not enough; A9/A10 price
discovery. Substituting a different router, softmax temperature, or fixed
pool would not move any cell. The claim under test is never ``\nm{} beats
architecture X''; it is that \emph{retention through dormancy requires an
assignment that is both reward-independent and covering, and
competition-then-pinning is a cheap emergent source of exactly that
conjunction}---with the objective axis an orthogonal dial governing
probe-window robustness to episode heterogeneity, not retention.
\end{keypoint}

\begin{remark}[Where the empirical record currently stands]
The companion experiments confirm the \nm{} cell directly: 20/20 seeds reach
a covering one-per-regime assignment; common regimes pin at
$\alpha \ge 0.95$; crisis acquires a stable unique owner at
$\alpha \in [0.67, 0.87]$ without a formal pin; never-pinning costs $0.0320$
vs.\ $0.0307$; $\lambda = 0$ costs $0.0314$ vs.\ $0.0307$; and the $W_c$
sweep is flat at $0.0303$--$0.0309$ with $W_c = 1$ marginally best---the
latter refuting the pre-registered batching prediction of
Section~\ref{sec:dd2-batching}. The refutation is reported with the same
prominence as the confirmations; a design-space map is only as trustworthy
as its falsified cells.
\end{remark}
